\section{Applications and discussion}
\label{sec:applications}

We close the paper with three short discussions. The first
(Section~\ref{ssec:non_trivial}) anticipates a referee objection
that Theorem~\ref{thm:empirical_invariance} is ``just CRT'' and
explains why the conclusion is not, in fact, an immediate
consequence of the tensor decomposition. The second
(Section~\ref{ssec:bell_implications}) summarises the consequence
for Bell-type non-locality, deferring the detailed discussion to
the companion paper~\cite{PT_HOLONOMY_NONLOCALITY}. The third
(Sections~\ref{ssec:open_k} and~\ref{ssec:future}) lists open
questions raised by the results, with the most subtle one ---
the closed-form expression for the depth parameter
$k_{p,q}^{\mathrm{eff}}$ --- already flagged in
Remark~\ref{rmk:k_interp}.

\subsection{Why the theorem is not trivial}
\label{ssec:non_trivial}

The CRT factorisation $T_m = \bigotimes_p T_p$
(Theorem~\ref{thm:crt_tensor}) and its corollary
$\pi_m^{\mathrm{Ruelle}} = \bigotimes_p \pi_p^{\mathrm{Ruelle}}$
(Lemma~\ref{lem:pi_factor}) are algebraic facts that depend only
on the structure of the residue ring $\mathbb{Z}/m\mathbb{Z}$ and
the spectral gap of each $T_p$. One might therefore suspect that
Theorem~\ref{thm:empirical_invariance} is a direct corollary of
these algebraic facts: since the joint law factorises, the
inter-channel mutual information vanishes, end of story.

This reading is correct under the Ruelle measure but incorrect
under the empirical measure. The two theorems together pin down
exactly where the algebraic argument is sufficient and where it
needs additional dynamical input. We organise the response under
four headings.

\paragraph{($\alpha$) The ``trivial'' identity $T_m \perp \mu$
hides a non-trivial choice of observables.}
The conclusion that $G(\mu)$ does not enter the empirical
inter-channel mutual information is not a statement about the
transfer matrix itself --- which is, by construction, independent
of $\mu$ --- but about the modular $\sigma$-algebras
$\mathcal{A}_p$ on which the observables $\Phi_p$ are defined.
A different choice of observables --- for instance, spatial fields
defined on the emergent metric $G(\mu)$ --- would lead to a
$\mu$-dependent conditional mutual information, even with the same
underlying gap field. The substantive content of
Theorem~\ref{thm:empirical_invariance} is the assertion that the
modular residues $r_n^{(p)}$ are the correct observables for the
question at hand, and that for this choice the $\mu$-dependence
collapses.

\paragraph{($\beta$) The closed form makes numerical predictions.}
Whether or not the proof of
Theorem~\ref{thm:empirical_invariance} is ``deep'' as a matter of
ergodic theory, Corollary~\ref{cor:closed_form} extracts from it
an explicit functional form for $\mathcal{I}_{p,q}$ that is fitted
by a single integer parameter $k_{p,q}$ per pair. The corpus
values $\mathcal{I}_{3,5} = 0.138$ bits and
$\mathcal{I}_{3,7} = 0.283$ bits are reproduced respectively at
$4.5\%$ and $8.1\%$ relative error, both within the
$\pm 20\%$ tolerance of the corpus T1 test. A predictive
closed-form expression with a single integer-valued fit parameter
is a non-trivial empirical assertion, regardless of the algebraic
elegance of the underlying proof.

\paragraph{($\gamma$) Without
Theorem~\ref{thm:empirical_invariance}, PT would be inconsistent
with Bell experiments.}
The empirical correlations between space-like separated detectors
in Bell-type experiments~\cite{Bell1964,AspectGrangier1982,Hensen2015}
are by now established at the eight-sigma level, with no
detector-separation-dependent signal of the kind that a
metric-mediated correlation would induce. If the inter-channel
mutual information $\mathcal{I}_{p,q}$ depended on the emergent
geometry $G(\mu)$, the theory would predict a measurable shift in
the experimental correlation as a function of detector separation
(which fixes the local value of $\mu$ in the Bianchi-I family).
Theorem~\ref{thm:empirical_invariance} rules out this signal at
the level of the theory itself, and is therefore the obstruction
that prevents Persistence Theory from being immediately falsified
by Bell-type experiments.

\paragraph{($\delta$) The proof is algebraically light but
conceptually heavy.}
At the level of mathematical content, Steps~1--4 of the proof of
Theorem~\ref{thm:empirical_invariance} use only the trivial
identity $I(X; Y \mid Z) = I(X; Y)$ for almost-surely constant
$Z$, the embedding $\sigma(\mathrm{Spec}(T_m)) \subseteq
\mathcal{I}$, and the Birkhoff--Hopf theorem. None of these
ingredients is technically deep. The non-trivial work is in
identifying the right division of $G(\mu)$ into a spectral object
on the one hand and the empirical measure $\pi_m^{\mathrm{emp}}$
into a trajectory object on the other, in such a way that the
combination of the two falls outside the joint information content
of the modular fibres. This division is the principal conceptual
content of the paper; the proof itself is a short verification
that the division survives.

\subsection{Implications for Bell-type non-locality}
\label{ssec:bell_implications}

The companion paper~\cite{PT_HOLONOMY_NONLOCALITY} provides a
foundational reading of Theorem~\ref{thm:empirical_invariance}.
The reading is the following: the modular fibres of the prime
sieve are not embedded into the emergent Bianchi-I geometry, and
therefore the inter-channel correlations they carry are not
correlations across space-like separated regions of an emergent
space-time --- they are correlations within a single arithmetic
substrate that happens to project, under the geometric flow, onto
space-like separated regions. The detector correlations measured in
Bell experiments are then artefacts of the projection, not signals
transmitted at superluminal speeds across the emergent
geometry. The result is therefore an existence proof, at the level
of measure theory and information theory, that a strictly local
realist substrate (the gap field $\{g_n\}$) can support the
empirical pattern of Bell-type correlations without contradicting
the no-signalling principle of the emergent geometry. We refer to
the companion paper for the detailed foundational
argument~\cite[\S~3--\S~5]{PT_HOLONOMY_NONLOCALITY} and for the
explicit connection to the Aharonov--Bohm
analogy~\cite{AharonovBohm1959}.

A direct consequence of this reading is that Persistence Theory
predicts no detector-separation-dependent
correction to the Bell correlations, contrary to the prediction
that a hypothetical metric-mediated extension would make. This is
a falsifiable prediction at the present experimental precision,
and is consistent with the most recent loophole-free Bell tests of
Hensen et al.~\cite{Hensen2015}.

\subsection{Open question: closed form for $k^{\mathrm{eff}}$}
\label{ssec:open_k}

The principal open question raised by the analysis is the
derivation of the empirical depth parameter $k_{p,q}^{\mathrm{eff}}$
from first principles. As discussed in
Remark~\ref{rmk:k_interp}, the fitted values
$k_{3,5}^{\mathrm{eff}} = 7$ and $k_{3,7}^{\mathrm{eff}} = 4$
differ from the CRT-naive depths $k_{3,5}^{\mathrm{CRT}} = 2$ and
$k_{3,7}^{\mathrm{CRT}} = 3$ in a way consistent with the
hypothesis that $k^{\mathrm{eff}}$ counts transfer-matrix
iterations through the antidiagonal double cover of $T_3$, but the
exact combinatorial formula remains conjectural. A derivation of
$k_{p,q}^{\mathrm{eff}}$ as a structural invariant of the pair
$(T_p, T_q)$, presumably from a finer spectral analysis of
$T_{pq} = T_p \otimes T_q$ on the active subspace, would close the
last remaining gap in the closed-form expression of
Corollary~\ref{cor:closed_form}. We leave this for future work.

\subsection{Formal verification status}
\label{ssec:formal_verification}

All algebraic results of the paper are formally verified in Lean~4
+ Mathlib under the namespace \texttt{PT.CrtDecoupling}, with
\textbf{0~\texttt{sorry}} and \textbf{0~\texttt{axiom}} across all
seven files of the verification (total 1246 lines of Lean source).
The verification is organised in five progressively stronger phases
that we summarise here for the reader.

\paragraph{Phase~1 (algebraic core).}
The CRT tensor factorisation
(Theorem~\ref{thm:crt_tensor}), the Ruelle factorisation
(Lemma~\ref{lem:pi_factor}), and the idealised Geometric Decoupling
Theorem (Theorem~\ref{thm:ruelle_zero}) are formally verified
without any measure-theoretic machinery:
\texttt{Tensor.crt\_tensor\_factor}, \texttt{Tensor.ruelle\_factor},
\texttt{Main.geometric\_decoupling\_ruelle}.

\paragraph{Phase~2 (conditional Theorem~6.1).}
The algebraic core of the conditioning step
(``conditioning on a $\pi$-a.s.\ constant random variable is a
no-op'') is formalised as
\texttt{Empirical.mutualInformation\_of\_a\_s\_constant\_factor}, together
with the analytic concentration lemma
\texttt{Empirical.expectation\_of\_constant}. Theorem~\ref{thm:empirical_invariance}
follows as \texttt{Empirical.empirical\_invariance}, conditional on
the a.s.\ constancy hypothesis on each $G(\mu)$.

\paragraph{Phase~2.5 (finite-state closure).}
The a.s.\ constancy hypothesis is discharged in the discrete
finite-state setting via spectral measurability of $G(\mu)$
(Lemma~\ref{lem:G_measurable}):
\texttt{SpectralReduction.geometric\_spectrally\_measurable},
\texttt{SpectralReduction.empirical\_invariance\_closed}. In the
finite setting, the $\sigma$-algebra $\sigma(\mathrm{Spec}(T_m))$
reduces to constant functions, so the a.s.\ constancy is immediate
without invoking Birkhoff--Hopf.

\paragraph{Phase~3 (infinite-state strengthening).}
For the general infinite path space, the result is obtained via
Mathlib's formalised Birkhoff--Hopf theorem
\texttt{Ergodic.ae\_eq\_const\_of\_ae\_eq\_comp$_0$}, packaged as
\texttt{Phase3.geometric\_a\_s\_constant\_ergodic} and
\texttt{Phase3.empirical\_invariance\_infinite}.

\paragraph{Phase~4 (sieve dynamical system).}
The concrete sieve dynamical system is modelled in
\texttt{Phase4.pathSpace}, \texttt{Phase4.shift},
\texttt{Phase4.shift\_measurable}, and packaged in the structure
\texttt{Phase4.SieveErgodicSetup} (state space, path measure,
ergodicity hypothesis). The closure theorem
\texttt{Phase4.sieve\_empirical\_invariance} applies Phase~3 to any
populated instance of this structure.

\paragraph{Phase~4.1 (chain closure: existence of a populated instance).}
A concrete populated instance \texttt{Phase4.trivialSetup} is
exhibited (state space $\mathsf{Unit}$, Dirac measure), with the
ergodicity \texttt{Phase4.shift\_ergodic\_unit} proved by direct
case analysis on the subsingleton path space. This proves
\texttt{Phase4.SieveErgodicSetup} is non-empty and makes
\texttt{trivialSetup\_empirical\_invariance} a non-vacuous
consequence of \texttt{sieve\_empirical\_invariance}.

\paragraph{Remaining technical gap (quantitative content).}
To obtain the empirical content of Theorem~\ref{thm:empirical_invariance}
for the \emph{actual} prime-sieve Markov measure $\pi_m^{\rm emp}$
(rather than the trivial Dirac instance), one must populate
\texttt{Phase4.SieveErgodicSetup} with a non-trivial instance built
from (i) \texttt{Mathlib.Probability.Kernel.IonescuTulcea} for the
Markov measure construction from the transfer-matrix kernel $T_p$,
(ii) the corpus primitivity of $T_p$ established in
\texttt{PT.Stochastic.T30FullSpectralAnalysis}, and (iii) the
classical theorem ``primitive Markov chain $\Rightarrow$ ergodic''
of ergodic theory --- not yet formalised in Mathlib at the level of
finite state spaces. This is the single remaining technical gap and
is a natural Mathlib pull-request.

The closed-form Corollary~\ref{cor:closed_form} is not formalised
as a theorem because the depth parameter $k_{p,q}^{\mathrm{eff}}$
is empirically fitted rather than derived
(Remark~\ref{rmk:k_interp}); a Lean-verifiable computational
statement of the closed form at a fixed $k$ would still be
straightforward and is documented in the verification README.

\paragraph{Reproducibility.}
The Lean files live in the project repository under
\texttt{PT\_LEAN/PT/CrtDecoupling/} and build with
\texttt{lake build PT.CrtDecoupling}. Total build time on a
modern laptop is approximately 15--30 seconds for incremental
builds of this subdirectory (Mathlib precompiled).

\subsection{Future directions}
\label{ssec:future}

Three natural extensions of the present work suggest themselves.

\paragraph{(a) Higher-order MI between three or more channels.}
The pairwise inter-channel mutual information
$\mathcal{I}_{p,q}$ does not exhaust the joint information content
of the modular fibres. The total correlation
$\mathrm{TC}\bigl(r^{(3)}, r^{(5)}, r^{(7)}\bigr)$ between all
three active channels, and its decomposition into pairwise and
higher-order components in the sense of
McGill~\cite[\S~7.4]{CoverThomas2006}, would refine the picture and
provide a stronger test of the closed
form~\eqref{eq:closed_form}. The corpus value
$\mathrm{MI}(r^{(3)}, r^{(5)}, r^{(7)}) \sim 1$ nat in the
aggregate suggests substantial higher-order content beyond the
pairwise $0.138 + 0.283 \approx 0.42$ bits.

\paragraph{(b) Extension to inactive primes $p \geq 11$.}
The echo primes $p \geq 11$ have anomalous dimensions
$\gamp{p}(\mustar) < 1/2$ and therefore do not contribute spatial
directions to the emergent geometry $G(\mu)$. They do, however,
contribute through the ghost vacuum polarisation mechanism of
chapter~22 of the monograph~\cite{PT_Monograph}, and the modular
$\sigma$-algebras $\mathcal{A}_{11}, \mathcal{A}_{13}, \ldots$
are non-trivial. An extension of
Theorem~\ref{thm:empirical_invariance} to include the echo
channels would clarify whether the empirical mutual information
is concentrated on the three active primes or extends into the
echo regime, and would constrain the geometric reading of the
ghost fraction $F_{\mathrm{ghost}}$.

\paragraph{(c) Connection to ergodic-theoretic invariants beyond MI.}
Mutual information is one entry in the family of information-theoretic
invariants attached to the shift dynamical system on the path
space. Other invariants ---
Kolmogorov--Sinai entropy~\cite{Walters1982},
mean information~\cite{KrengelBook}, and the higher
correlations of the transfer operator~\cite{Seneta2006} --- all
admit empirical-versus-Ruelle splittings analogous to the one
studied here. A systematic investigation of which of these
invariants are geometric-invariant in the sense of
Theorem~\ref{thm:empirical_invariance} would clarify the structural
content of the result and place it in a broader ergodic-theoretic
context.
